\section{Experiments}

\subsection{Mate-in-1 Probing}
\paragraph{Research question} Do maia networks know they play bad, like do they know when they are doing blunders and do they know the best move?

\paragraph{Method} We look at a mate in 1 dataset. On this dataset, we perform a linear probing on the model. We probe the model on the binary classification task (is it mate in 1?) and the exact move task (what is the best move?). We use double linear probing (predicting the start / end squares) to decompose the move prediction task into two separate tasks.
The evaluation should also be done on a randomised model.

\paragraph{Results} The binary classification is clear, the model know when a position is a mate in 1 even on the failed puzzles. The results also indicate that the move can be predicted confidently.

\subsubsection{Localised Probing}
\paragraph{Research question}
Can we localise precisely where the best move knowledge is located?
\paragraph{Method}
train one probe per layer, maybe also only use one square as input for the probe, maybe average / maxpool over the squares.

\subsection{Value Decomposition}
\paragraph{Research question} Can we decompose the value into piece value and position (centrality or activeness)? For different stage of the game? In terms of controlled squares?

\subsection{Principal Variation}
\paragraph{Research question}
Same first best move but different sequence. Or same puzzle but different sequence. Should not be able to predict the line of play agains a random model or a far stronger model.
\paragraph{Method}
Simulate self play with MCTS or distributional sampling.
Predicting multiple moves in parallel. Recover from the fact that moves are maybe not considered at all the layers.

\subsection{Summed Explanations}
\paragraph{Research question}
Summed explanations for MCTS explanations
\paragraph{Method}
Evaluate with counterfactuals on the MCTS tree.
Train a model to predict the output given the saliency map.
Contrastive counterfactuals for evaluating saliency maps

\subsection{Latent Saliency Maps}
\paragraph{Research question}
Explore the saliency maps of the latent representation.
\paragraph{Method}
Latent saliency patching for saliency map evaluation -> quantus
Same as above.


Impact of history
